情感由语言内容、语音韵律与面部动作共同表达，但三路信号在采样粒度与时间尺度上存在差异，真实场景还伴随局部信息缺失与判断依据难以追溯等问题。本文基于赛题提供的CMU-MOSEI数据，分别建立词级时间锚点对齐、局部缺失鲁棒预测与输入干预解释模型，并以统一的时间组织与掩码表示衔接三项任务，实现原始素材的特征组织、缺失条件下的情感预测与关键证据定位。

针对问题一，采用CTC强制对齐估计原词时间区间，由字符偏移建立子词与原词的归属关系，并按时间交叠长度聚合语音帧与视频帧，建立不确定边界锚点对齐模型。完成附件1全部100条视频的特征提取，文本、语音和视觉特征维度分别为$50\times768$、$50\times25$和$50\times16$；1932个原词中1920个获得有效时间区间，生成率为99.38\%。通过边界核验、扰动敏感性与典型样本分析验证时间锚点的有效性，生成的特征、掩码与时间记录均可回映至原始素材。

针对问题二，建立三态感知可靠补全与融合模型TSRF-D，以结构位、可观测位和缺失位刻画局部缺失状态，通过样本内候选补全与可靠度门控融合上下文，并利用冻结教师模型的分类与强度软监督训练学生模型。在105种连续局部缺失条件下，模型平均Macro-F1为0.6066、MAE为0.6492，在比较模型中取得最高Macro-F1和最低MAE。进一步分析表明，缺失模态类型和缺失程度是性能退化的主要因素：文本缺失造成的影响明显高于语音和视觉，且连续缺口增大时预测性能持续下降，而缺失位置的影响相对较小。模型完成附件3全部30条样本的极性与强度预测。

针对问题三，构建ICF-Shapley解释模型，通过枚举模态联盟分解文本、语音和视觉对分类与回归输出的贡献，并利用局部删除响应定位关键片段。结果表明，文本是总体主要参考模态，但部分样本仍表现出明显的语音或视觉主导特征；在冻结模型的测试子集上，高贡献片段的删除对分类和强度输出造成的影响均显著高于随机片段，表明所定位证据能够反映模型实际依赖的输入信息。完成附件4全部20条样本的预测与解释，118条可用证据均成功回映至原始文本、语音时段和视频关键帧，同时发现少数非文本主导归因对解释参考基线较为敏感。

\keywords{多模态情感分析、时序对齐、局部模态缺失、教师--学生蒸馏、Shapley解释}
